%-----------------------------------------------------------------------------%
\chapter{\babEmpat}
\label{bab:4}
%-----------------------------------------------------------------------------%
This chapter describes the method used to evaluate the effectiveness of Latviz as a visualization tool for ML-KEM, ML-DSA, LLL, and BKZ as well as the results of the evaluation.

\section{Methods}
 To evaluate the effectiveness of Latviz as a visualization tool for ML-KEM, ML-DSA, LLL, and BKZ, a survey of Latviz users was conducted. The primary goal of the survey was to determine whether or not Latviz helped users understand the four algorithms, how effective it was at doing so, and which parts of the tool were particularly effective.
 The group of correspondents consisted of ... students from the computer science faculty of the University of Indonesia. The survey was distributed directly to students within the researchers' academic network, focusing primarily on students who had taken cryptography.
 The survey consists mostly of questions resembling those of SALG (Student Assessment of their Learning Gains), which is a survey designed to measure students' gains in specific knowledge, skills, and attitudes \cite{selfrep}. The questions follow along the lines of "As a result of using Latviz, what gains did you make in..." and "How much did this aspect of Latviz help in...", which measure knowledge gain and how much of it can be attributed to specific parts of Latviz respectively. Additionally, an optional section was provided to gather user suggestions or feedback that was not covered by the available questions.
 Each algorithm also had pre- and post- test sections for more objectively measuring improvements in a user's knowledge. Improvements are measured using normalized gain \cite{PhysRevPhysEducRes}, which is the fraction of knowledge learned that was not known prior to taking a course, or in this case, prior to using Latviz. It is formally defined as:
\[
g = \frac{P_{\text{post}} - P_{\text{pre}}}
{100\% - P_{\text{pre}}}
\]
where $(P_{\mathrm{pre}})$ and $(P_{\mathrm{post}})$ denote the average pre-test and post-test scores respectively.
(Mungkin paragraf ini dimasukin ke results? Thoughts? abisnya dipake untuk analisis, jadi i rasa mungkin bisa dimasukin results aj?)

\section{Results}
On the whole, the results of the user survey suggest that Latviz did improve users' understanding of ML-KEM, ML-DSA, LLL, and BKZ, and was also considered an effective tool for learning the algorithms. 
